\section{Turning the Tables}

We have given free rein to Guenon's critique of Evola which, as far as it goes, is justified. Now we can turn the tables and explore what Evola found lacking in Guenon, specifically, the two issues of ``Guenonian Scholasticism'' and ``Bureaucratic Initiation''. The mutual critiques are of different orders, but that does not make one less important. The distinction is between ``thesis'' (the principle) and ``hypothesis'' (the application of the principle to the concrete situation).

Guenon is unremittingly concentrated on the thesis in his published writings, which are typically abstract and require some effort to parse. Even in works such as Crisis of the Modern World or Reign of Quantity, the concrete situation is seen in negative terms. Evola, on the other hand, seems impatient with that abstraction, and focuses on the hypothesis, i.e., dealing with one's concrete situation as he famously does in Ride the Tiger and Men among the Ruins.

\paragraph{Scholasticism}


Guenon's most metaphysical works (Multiple States of Being, Man and his Becoming, Symbolism of the Cross) may seem to many as too abstract and overly logical. It takes some time to integrate that system of thought into one's worldview and learn to experience the world from that vantage point, \emph{sub specie aeternitatis}, as it were. Yet, Guenon asserts it is above logic, while not illogical, and requires a ``realization'' to grasp it in its depths.

So without such a realization, it remains abstract and can devolve into a type of Scholasticism, as happened, for example, with Thomism. Examples can be found in some neo-Traditionalist writings, printed or on the Internet, which are often abstract discussion about whether such and such ancient religion was Traditional, or the same for some person or another, as though they are really in a position to discern that.

To avoid such a dry Scholasticism, Guenon informs us of the need for a real initiation from a valid organization. Unfortunately, Guenon is reluctant to share any specifics. He, for example, explains the principle of the Supreme Identity, but not how to realize it. He gives the Divine Comedy as an example of a initiatory text that outlines the path, but ignores any specific details. For example, he will point out that the stages of that poem represents higher states of being, but the reader is on his own to explore what that means as far as a concrete application.

\paragraph{Bureaucratic Initiation}


Guenon does seem to describe initiation in a bureaucratic way: first, a man must discover a valid organization and then follow their process in a precise and well-defined way. Yet this begs the question of how can a man, starting from a concrete situation of ignorance, be in a position to recognize what is valid or not?

Here, I suppose, it is helpful to know the metaphysics which can serve as the touchstone to evaluate a teaching. Nevertheless, there is no ISO\footnote{\url{http://www.iso.org/}} standard to rate such teachings. The other objection to such a strict view is that the ``Spirit blows where He wills'', so organizations may decline and arise apart from any human effort.

Now I get several newsletters about health and nutrition, which all contradict each other. There is the paleo diet, the Mediterranean, vegan, and so on, mutually opposed and often polemical against each other. On the one hand, there may be frustration about which is the ``ideal'' human diet, but another conclusion is that there are different diets suitable for different types of men.

Evola, to his credit, made an effort to identify initiatic paths suitable for a modern Westerner. In the Ur/Krur group, he interacted with Theosophists, Anthroposophists, Jungians, inter alia. At one point, he even spent some time in a Catholic monastery, if you can envision him as a monk. He eventually settled on several figures whom he regarded as offering something, among whom are: Eliphas Levi, Kremmerz, Meyrink, Bo Yin Ra, Gurdjieff, Crowley. Perhaps they are ``irregular'' from Guenon's point of view, yet they clearly understood something and learned it from somewhere. A comparative study of their teachings may be of interest, if the right person ever undertakes it.

Guenon makes the distinction between the lesser mysteries and the greater mysteries, and let us suppose for the sake of argument that Guenon's real objection is that those figures had an inadequate understanding of the greater mysteries, perhaps even because they were not aligned with one of the valid traditions. Nevertheless, knowledge of the lesser mysteries may very well be the ``proper diet'' for many men today.

While Guenon criticized Evola's interest in those men, he was not forthcoming about an alternative. In his published writings (as far as I recall), he denied the existence of initiatory organizations in the West, yet in his letters he claimed to have received both a Western and an Eastern initiation; he never revealed from where, as far as I know. The reason can only be speculated; perhaps, the group wished to remain unknown, yet somehow Guenon found them.

In my experience, I have met many people involved with various Vedanta inspired organizations, been to their meetings, heard their talks. Not to be judgmental, it is hard to discern much difference between them and the general populace. Belonging so such an organization often seems to be more of a social thing, a place to make friends, eat exotic foods, and establish an identity. Granted, there may be those with a deeper understanding, but the point I'm making is that initiation in and of itself is insufficient.

In East and West, Guenon passed over Tibetan Buddhism on the grounds that it is too remote from Westerner's understanding. But that was in the 1920s and today Tibetan Buddhist centers can be found in every major city in Europe and North America. Unfortunately, its encounter with the West has not brought the West to Tradition, but seems, on the contrary, to have affected Buddhism for the worse. But that is another topic; suffice it to say for now that the ideal of a theocratic, patriarchical, hierarchical society is not the ideal of any Western Buddhists.

\paragraph{Fraud and Legitimacy}


While it may be difficult to find a legitimate teaching, a fraudulent one may be easier to identify. The most obvious is if the organization is used for the personal gain of the leader. Hermetists did not earn their living through their teachings. Instead, they held positions, such a horse traders or street artists, which gave them the freedom, and the cover, to travel from town to town to meet with their students. The second is that whether the organization is able to hand down its tradition, which is its very nature, or whether it is just a way of life for its successors.

As an example, I can point to Valentin Tomberg, who explicitly requested that no organization be formed based on his work. This does arouse some curiosity, doesn't it? Why did he expect that to happen? Is doesn't seem likely that he was so vain he assumed it. Perhaps he realized its importance and suspected it would have an impact, even if not right away. Or even perhaps, he was part of an movement that decided this was the opportune time to leave behind some evidence; he certainly hints at this.

In the example of the \emph{Fedeli d'Amore} --- which includes Dante, the authors of the Romance of the Rose, Boccaccio --- the organization apparently no longer exists, yet it did make sure to leave written texts. As Guenon points out, it took several hundred years for anyone to realize that the Divine Comedy was an esoteric work and not just a poem.

For more on the notions of thesis and hypothesis, and accommodation of principles to concrete situations, interested readers can read this\footnote{\url{https://opuscula.blogspot.com/2009/08/intervention-of-mgr-dupanloup-part-i.html}}.


\flrightit{Posted on 2012-08-10 by Cologero}

\begin{center}* * *\end{center}

\begin{footnotesize}
\begin{sffamily}
\texttt{Avery Morrow on 2012-08-10 at 01:22 said:}

I actually was considering precisely the same problem at roughly the same as you were writing this post. This is the conclusion I came up with:

\begin{quotex}\footnotesize

The ``unity of tradition'' doesn't mean that religions all teach the same thing, but that the most developed of these paths, using its own language, leads to the same wisdom as the other well-developed paths.

Some projects, like Theosophy and the New Age, are anti-tradition and go around in circles forever. After talking with my friend about Bahai Faith, Tenrikyo, and Mormonism, which do seem to guarantee spiritual peace to adherents, I think I understand the meaning of ``illusory tradition'' as defined by Guenon here. Adherents of these new movements will enter into a journey. But their new paths are undeveloped, and adherents do not yet understand how to see the eternal nature of tradition and cannot receive initiation using the language-worlds created by these vehicles. The average person can live happily with imperfect understanding, so these ``illusory traditions'' are not useless as ordinary religions, but the intellectual will remained mired in modernity and will have to force his way to escape. Only an old tradition can really guarantee intellectual engagement.

The importance of finding an old tradition has grown for intellectuals because modernity, anti-tradition, builds a language-world that completely rejects the possibility of tradition, a rejection that did not exist in the past or was more widely disputed. This has created confusion and chaos when discussing issues of the mind, which has caused many of these ``illusory traditions'' to appear. But a real tradition will have a very strong vocabulary from which modernity can be observed and contrasted with truth.
\end{quotex}


I might add to this that Guenon and Evola may be addressing different problems. The authority that Guenon wants to retrieve from Tradition is a font of sacred knowledge. The hierarchy he proposes is a hierarchy of wisdom, in which knowledge of the profane world symbolizes more constant, eternal concepts which are grasped through initiation. Evola's main interest, to generalize quite roughly, is Tradition as a font of power that flows into the occult war. His World-King really is a world king and not a symbolic creature. Recognizing what is traditional here is recognizing the forms of traditional authority as exerted on the unintellectual world. Thus Evola is not a disciple of Guenon in the same way that the ecumenist Schuon is, because although he agrees on the existence of Tradition, he is asking different questions about it.

\hfill

\texttt{Izak on 2012-08-10 at 03:11 said:}

I thought that Guenon continued to proclaim the legitimacy of the Freemasons up until the day he died. My guess is that he was initiated with them.

In the pamphlet ``Guenon: Teacher for Modern Times,'' Evola takes him to task on his view of the Masons, and publishes a couple of Guenon's epistolary responses regarding the subject. I tend to agree more with Evola, and I also agree with Evola's point about the basic dearth of tradition in the East.

The problem with these so-called initiatory organizations is that they can only do so much as isolated entities. What I see as Evola's chief virtue is that he recognized that a traditional civilization must contain a uniting of both the spiritual and temporal order with which to give faith coherence. This is probably why he admirably never gave himself over entirely to (IE, was initiated by) a corrupt or flawed organization. He recognized that he would never achieve a total immersion, or that too much corruption would seep in, and he thus promoted an individualistic approach --- take it or leave it.

In one of Guenon's letters to Evola, I seem to remember him giving a very sparse few examples of initiatory organizations in the West, though he is unsure as to whether they count for anything, or whether they have merely degenerated into ``study groups'' (his term!).

I don't know anything about Evola's criticisms of Guenon's ``scholasticism,'' but it would seem somewhat pointless in my mind for Evola to even bother leveling such an accusation in the Kali Yuga, where just about any discussions of this nature are bound to become needlessly abstract and/or dialectical. It would probably be more appropriate to accuse Guenon of producing works which would inevitably engender a type of scholasticism.

In my view, the real problem with medieval scholasticism is that it basically neglected the eschatological/anagogical sense of scriptural exegesis --- the only sort which dealt with true metaphysics --- which early medieval tradition handled admirably (\url{http://www.ecriture-art.com/raban-e.htm\#} -- here's The Symbolism of the Cross, 9th century ed.). By the 12th century, a huge wave of inferior theology took over the medieval tradition with the likes of Abelard (while people like Bernard of Clairvaux and some of the Victorines stayed true). Whatever true metaphysics Thomas was able to resuscitate was overshadowed by the rigor of his form as opposed to content, and the nosedive continued undaunted. Calling Guenon a scholastic seems a bit silly to me, since metaphysics was his chief concern. It isn't as though Guenon was sitting around writing speculations as to whether or not the nephilim were the products of demon incubi impregnating mortal women during astrologically significant days on the calendar (as our friend Thomas, in fact, did!). All of his questions were significant.

\hfill

\texttt{Mihai on 2012-08-10 at 05:20 said:}

Guenon explicitly declared that he wished to explain certain traditional ideas and doctrines in a way accessible to the modern western mind- a mind that is especially unaccustomed to the direct perception of symbolism. Hence his more logical and dialectical approach.

As for the question of initiation, Guenon was initiated in a Sufi lineage, this is the only thing certain. As for the western initiation, it is really unknown to us.

Guenon stressed the importance of initiation in an orthodox tradition, regardless of which. I believe this to be legitimate, since the only alternative to this is some neo-spiritualism or an attempt to find spirituality as a freelancer, which I believe not only to be dangerous, but also delusional and with no chances of success.

I have read Evola's alternatives that he wrote in ``The limits of initiatory regularity''. He was certainly right to call Guenon on the fact that he does not offer any practical advice, and leaving the seeker in the dark as to the solution, but his ``alternatives'' aren't really that convincing, and in some ways quite fanciful.

He talks about dream initiation or a certain type of spontaneous event that would initiate the person without the help of an intermediary, but such are really extremely rare cases which, for all practical purposes, should not be relied upon by any person seeking initiation. Besides, such cases have always been rare and in the case of especially ``gifted'' people. It is easy to consider oneself, in one's own eyes, as being one of such ``gifted'', but who can tell one is not deluding oneself ?

Besides, as his life showed, Evola never did find the solution to this question. His works tend to show that. Of course, he gave valuable insights and glimpses of alternatives to modern materialism, that can be expanded upon by different people, but still, the fact that he wrote about such diverse topics- Tantra, Buddhism, Alchemy, mountain climbing, esoteric sexuality etc.- also show his inward seeking for a path of realization, a path that he never came to really grasp.

On the other hand, it is easy to see, throughout Guenon's writings, a continuity and a logical unfolding with each single work.

PS: I am not aware if Guenon was aware of Crowley in any way, but all I can say, regarding him, is that he held some valuable knowledge, but ultimately he gave birth to a dangerous and syncretism of half-truths, packaged in an evolutionist and individualist mentality. His ``Thelema'' is really indistinguishable from other modern spiritualisms and occultisms that announce a spiritual rebirth of humanity.

\hfill

\texttt{Gabe Ruth on 2012-08-10 at 13:13 said:}

This is off the topic of the interesting post, but Avery, your note about the nature of these new religions brought this to mind:

``For you, the world will continue to wear a noble, awful face. You will never rise above mysticism, but be happy in your own way.''

It's from A Voyage To Arcturus, by David Lindsey. I would be very interested to hear the thoughts of anyone here on it. I can't recommend it, but it is a startlingly interesting book.

\hfill

\texttt{Andrew on 2012-08-11 at 23:27 said:}

The point about Buddhism being negatively effected by its encounter with the West is quite right. Everything that has encountered modernity has been corroded by it. This puts me in an awkward position as I am a Buddhist at odds with nearly every other ``Buddhist'' in America. As I see it, the only option is to apply principles to practice and form friendships with like-minded people. With this, we build fortresses and wait out the storm. Though in line with ``Ride the Tiger'' there may be a place for those unafraid to make matters worse so as to collapse our current culture sooner.

\hfill

\texttt{Aghorable on 2012-08-15 at 01:30 said:}

Good article. ISO standards for teachings, that's comedy. I don't look forward to the day of *ISO 16003940 -- Initiation Delivery Systems*.

``... it is hard to discern much difference between them and the general populace.''

That observation might be explained by an important distinction; that between virtual initiation and actualisation. If my understanding is not flawed, this latter depends on a truly unique catalyst to get it started, one that cannot exactly be imposed from outside. But without actualisation, there is only a potentiality that has not yet blossomed, albeit a determinate one.

One can also ask what would be the purpose of receiving more than one initiation, if it is only a question of metaphysical realisation. On the other hand, if it is also about obtaining an alignment that would allow one to participate in a particular organisation's work, then it can be understood.


\hfill

\end{sffamily}
\end{footnotesize}

